\thispagestyle{otFront}
\phantomsection\label{references}\pdfbookmark[0]{Références}{bm-references}
\otPageTitre{Références}{France · Terminale générale · Spécialité physique-chimie}
\fontsize{10.5}{14.5}\selectfont
\otIntertitre{Le programme de référence}
Le livre s'appuie sur le programme de l'enseignement de spécialité de physique-chimie de la
classe terminale de la voie générale, publié au \textbf{Bulletin officiel spécial n° 8 du
25 juillet 2019} (NOR MENE1921249A) et en vigueur depuis la rentrée 2020. Ses quatre thèmes
structurent les parties du livre ; la partie « Mesure et incertitudes » du programme fait
l'objet du chapitre 27.\par\vspace{3mm}
\href{https://www.education.gouv.fr/bo/19/Special8/MENE1921249A.htm}{\textcolor{otNavy}{Programme officiel sur education.gouv.fr}}\par\vspace{5mm}
\otIntertitre{L'épreuve du baccalauréat}
La préparation suit la \textbf{note de service du 11 septembre 2026}, NOR MENE2622644N,
publiée au Bulletin officiel spécial n° 4 du 17 septembre 2026. Applicable à compter de la
session 2027, elle remplace la note de service du 11 février 2020. L'épreuve porte sur
l'ensemble du programme de terminale ; les notions de première non approfondies en terminale
doivent être connues et mobilisables, sans constituer un ressort essentiel du sujet.
La partie écrite (3 h 30, trois exercices) consacre deux points sur vingt à la maîtrise de la
langue. La partie pratique (1 heure) évalue les compétences expérimentales.\par\vspace{3mm}
\href{https://www.education.gouv.fr/bo/2026/Special4/MENE2622644N}{\textcolor{otNavy}{Définition officielle de l'épreuve de spécialité physique-chimie}}\par\vspace{5mm}
\otIntertitre{Les choix pédagogiques}
Les exercices, les défis et les sujets type bac sont des créations pour ce livre ; ils ne
reproduisent aucun sujet officiel. Les défis très exigeants sont facultatifs et indiquent
leurs prérequis ; les outils qui dépasseraient le programme sont fournis dans l'énoncé.
Les données réelles sont accompagnées de leur source, consultée à la date de l'édition.\par\vspace{3mm}
Les schémas, courbes et tableaux sont vectoriels. Les programmes Python peuvent être exécutés
séparément. Les résultats numériques sont donnés avec un nombre de chiffres significatifs
cohérent avec les données.\par\vspace{5mm}
{\fontsize{9}{12}\selectfont\color{otMuted}Sources officielles consultées le 21 octobre 2026.
La charte reprend le modèle OptiTAB : couleurs, typographie et repères de difficulté.}
\clearpage
